\begin{frame}{Local Optimization}
\framesubtitle{Choosing Retained Tensors}

\begin{itemize}
  \item 그룹의 텐서를 느린 메모리를 거치지 않도록 지정할 수 있습니다.
  \begin{itemize}
    \item 빠른 메모리 공간을 차지하기 때문에, 꼭 필요한 텐서만을 선택해야 합니다.
  \end{itemize}
  \item Operator fusion과 같이, 무식하게 모든 경우를 탐색할 수는 없습니다.
\end{itemize}
\pause
\begin{figure}
  \centering
  \only<1-2>{\includegraphviz[width=0.8\linewidth]{graphs/bad-retention}}%
  \only<3->{\includegraphviz[width=0.8\linewidth]{graphs/good-retention}}
\end{figure}
\pause
\begin{tcolorbox}[title={제약 조건}]
  빠른 메모리에 유지할 텐서는 그룹의 경계에 있어야 의미가 있습니다.
\end{tcolorbox}
\end{frame}

\begin{frame}{Local Optimization}
\framesubtitle{Partitioning Graphs}

\begin{itemize}
  \item 추가적으로, 유지되는 텐서는 최대 한 개의 그룹에서만 사용 가능합니다.
  \item 이 관점에서, 유지되는 텐서를 선택하는 것은 선택된 그룹의 분할 문제와 같습니다.
\end{itemize}
\pause
\begin{figure}
  \centering
  \only<1-2>{\resizebox{!}{2.5cm}{\includegraphviz[width=\linewidth]{graphs/before-retention}}}%
  \only<3->{\resizebox{!}{2.5cm}{\includegraphviz[width=\linewidth]{graphs/after-retention}}}
\end{figure}
\pause
\begin{itemize}
  \item 실행 그룹의 각 텐서들에 대해 유지 여부를 결정해 추가적인 최적화를 구현했습니다.
\end{itemize}
\end{frame}